\documentclass[conference]{IEEEtran}
\usepackage{amsmath,booktabs,url}
\title{Phase 2 Progress in Learned Reverse Tracking and Docking}
\author{
  \IEEEauthorblockN{Nandan Prince}
  \IEEEauthorblockA{MSML 642}
  \and
  \IEEEauthorblockN{Chris Pedretti}
  \IEEEauthorblockA{MSML 642}
}

\begin{document}
\maketitle

\begin{abstract}
We evaluate initial PPO training for articulated reverse docking in the 2D
kinematic simulator. The open-bay policy learned to sustain reverse motion
from an aligned start and occasionally docked during training. Its final
deterministic rollout reached the target vicinity and missed the heading and
trailer containment requirements.
\end{abstract}

\section{Setup}
The compact tractor--semitrailer reverses at a target speed of
$0.20\,\mathrm{m/s}$ with normalized steering rate as the PPO action. The
open bay has no physical dock walls. The painted target, articulation and
workspace limits, and nominal docking criteria still apply. Docking requires
longitudinal and lateral errors of at most $0.10\,\mathrm{m}$, heading error
of at most $3^\circ$, and full trailer containment. The aligned-start
experiment used seed 0, reward version 1, and a budget of 200,000 steps.

\section{Aligned-Start PPO}
The run completed 200,704 environment steps and 416 training episodes.
There were 10 successes, with the first at step 79,748.
Table~\ref{tab:training} shows the evolution of episode outcomes.
Training sampled steering actions; evaluation used deterministic actions.

\begin{table}[ht]
\caption{Selected Training Windows by Episode Completion Step}
\label{tab:training}
\centering
\begin{tabular}{lrrr}
\toprule
Window & Jackknife & Workspace exit & Success \\
\midrule
First 20,000 steps & 80 & 1 & 0 \\
After step 150,000 & 0 & 77 & 8 \\
\bottomrule
\end{tabular}
\end{table}

The final policy exited the workspace in all 20 repetitions of the same
aligned start. At the target crossing, its lateral error was
$0.088\,\mathrm{m}$ and absolute heading error was $5.41^\circ$. The full
trailer footprint was outside the containment region. Peak articulation
was $8.38^\circ$ in the final replay, compared with $45.36^\circ$ at the
20,000-step checkpoint. The 80,000-step checkpoint was fully contained at
the target crossing, with $0.099\,\mathrm{m}$ lateral error and
$3.29^\circ$ heading error. Zero steering docked in 486 steps.

\section{Centerline Recovery}
Recovery uses the same open bay and reward with initial lateral offsets
sampled up to $0.15\,\mathrm{m}$ in either direction and aligned
tractor--trailer headings up to $2^\circ$ from the dock axis. The reference
stays on the dock centerline for every reset. Both recovery runs started from
the final aligned policy weights with fresh PPO optimizers.

After 100,352 additional steps, training recorded 126 workspace exits and
48 jackknifes across 174 episodes. The final policy exited the workspace
on all 20 perturbed validation starts, missing the 70\% advancement gate.
On the same 20 starts, the initialization policy also docked zero times
but finished with a mean absolute lateral error of $0.31\,\mathrm{m}$,
compared with $4.56\,\mathrm{m}$ after further training. Zero steering
docked in 7 of 20 cases. The lateral errors are measured at termination.
These observations show worse terminal alignment for this warm start
after training on the perturbed distribution.

A second recovery run with the longer budget completed 200,704 steps and
347 episodes: 2 docks, 68 jackknifes, and 277 workspace exits. The final
deterministic policy jackknifed on all 20 validation starts. This run
did not produce reliable recovery despite the longer budget.

\section{Gentle-Curve Feasibility Check}
The next PPO experiment uses a single smooth reference with the trailer tail
starting $8.0\,\mathrm{m}$ beyond the bay opening and $0.50\,\mathrm{m}$
off-center, with zero initial heading offset and articulation. The bay remains
open. In one deterministic setup check, LQR docked in 687 steps within the
900-step limit, and zero steering exited the workspace. PPO training on
this specific curve is pending. A later paired comparison on identical starts
and goals will test the effect of path-relative guidance.

\section{Evidence and Limitations}
The aligned-run logs are stored under campaign
\path{20261125-2331_ppo_ppo-ui-campaign}. The single-seed evidence
characterizes sustained reversing and the observed docking failures.
Recovery provenance and the final per-episode metrics are saved under
\path{20261125-2346_ppo-curriculum_ppo-open-dock-recovery-v1}.
The longer recovery campaign is
\path{20261125-2359_ppo_ppo-ui-campaign}.
Causal effects of removing the dock walls and performance relative to NMPC
remain open questions.

\end{document}
